\documentclass[11pt]{article}
\usepackage[margin=1in]{geometry}
\usepackage{booktabs}
\usepackage{amsmath}
\usepackage{amssymb}
\usepackage{hyperref}
\usepackage{cite}

\title{Systemic Pharmacological Suppression of Neural Activity Reverses Learning Impairment in a Mouse Model of Fragile X Syndrome}
\author{}
\date{}

\begin{document}
\maketitle

\begin{abstract}
The enhancement of associative synaptic plasticity often results in impaired rather than enhanced learning. We previously proposed that such impairments reflect a history-dependent change in the threshold for plasticity. We extend this work using a second mouse line with enhanced parallel-fiber (PF) to Purkinje cell long-term depression (LTD), the Purkinje-cell-specific Fmr1 knockout (L7-Fmr1 KO) model of Fragile X syndrome (FXS). L7-Fmr1 KO mice were selectively impaired on two LTD-dependent oculomotor learning tasks: high-frequency (1 Hz) VOR-increase learning (WT 0 vs 30 min: $p = 7.37\times10^{-4}$; L7-Fmr1 KO 0 vs 30 min: $p = 0.97$; L7-Fmr1 KO vs WT at 30 min: $p = 3.47\times10^{-5}$; Tukey) and OKR adaptation ($p = 1.27\times10^{-4}$ L7-Fmr1 KO vs WT at 60 min; $p = 0.87$ L7-Fmr1 KO 0 vs 60 min). VOR-decrease learning was normal ($p = 1.10\times10^{-5}$ L7-Fmr1 KO 0 vs 30 min; $p = 0.091$ L7-Fmr1 KO vs WT at 30 min). VOR-decrease pre-training and Vestibular-only pre-training each rescued the VOR-increase deficit ($p = 0.001$ L7-Fmr1 KO 0 vs 30 min after VOR-decrease pre-training; $p = 0.84$ L7-Fmr1 KO vs WT after Vestibular-only pre-training). Pre-treatment with diazepam (0.5 mg/kg i.p.) 18--24 hours before training reversed the high-frequency VOR-increase deficit ($p = 0.86$ L7-Fmr1 KO vs WT at 30 min) and the OKR deficit ($p = 0.02$ L7-Fmr1 KO vs WT at 60 min, in the opposite direction) in L7-Fmr1 KO mice without affecting baseline VOR gain or VOR-decrease learning. The learning-enhancement effect of diazepam was transient: tested 1 week later, L7-Fmr1 KO mice again failed to learn ($p = 0.12$). Low-frequency (0.5 Hz) VOR-increase impairment was not rescued by either pre-training or diazepam. These results support a sliding-threshold metaplasticity hypothesis and suggest a clinical avenue for learning disorders characterized by enhanced associative plasticity.
\end{abstract}

\section{Introduction}
Long-term synaptic plasticity has long been a therapeutic target for brain disorders affecting learning and memory. Some manipulations that augment plasticity do augment learning \cite{tang1999genetic,vanpraag1999running,lee2009learning}, but many impair learning \cite{migaud1998enhanced,uetani2000impaired,gu2002protein,cox2003visual,rutten2008enhanced,navakkode2022parvalbumin}. Nguyen-Vu et al. \cite{nguyenvu2017saturation} proposed that manipulations enhancing PF-Purkinje cell LTD can impair learning by allowing spontaneous activity to aberrantly recruit LTD, increasing the threshold for subsequent induction. This was supported in MHCI K$^b$D$^b-/-$ mice with enhanced PF-Purkinje cell LTD \cite{mcconnell2009mhci}: optogenetic LTD induction in WT phenocopied the knockout's VOR learning deficit, and VOR-decrease pre-training reversed the deficit.

Here we test this hypothesis in a second mouse line with enhanced LTD, L7-Fmr1 KO (Purkinje-cell-specific Fmr1 knockout, Fragile X model) \cite{koekkoek2005deletion,mientjes2006conditional,zhang2004l7cre}. We further test a new prediction: pharmacologically suppressing neural activity in the pre-training period should prevent LTD induction and restore subsequent LTD-dependent learning capacity.

\section{Methods}
\subsection{Animals}
L7-Fmr1 KO mice were generated by crossing conditional Fmr1 knockout \cite{mientjes2006conditional} with L7/Pcp2-Cre \cite{zhang2004l7cre}. WT littermates served as controls. Both sexes were tested.

\subsection{Oculomotor learning tasks}
VOR was measured with 1 Hz sinusoidal rotation about an earth-vertical axis with peak velocity $\pm 10^\circ$/s. VOR-increase training paired vestibular stimulation with oppositely directed visual motion for 30 min. VOR-decrease training paired vestibular stimulation with same-direction visual motion. OKR adaptation was induced with optokinetic drum rotation (1 Hz, peak velocity $\pm 10^\circ$/s) for 60 min. Low-frequency tests used 0.5 Hz. Vestibular-only pre-training used vestibular stimulation alone in darkness.

\subsection{Diazepam administration}
Diazepam (0.5 mg/kg, or 0.4--0.5 mg/kg, or 2.5 mg/kg i.p.) was administered 2 hours or 18--24 hours before training in separate cohorts. Saline was used as a control. The 18--24 h delay exploited diazepam's approximately 24-hour half-life \cite{riss2008benzodiazepines}.

\subsection{Statistics}
One- and two-way ANOVA with Tukey post-hoc or two-sample/paired-sample t-tests were used. All plotted values are mean $\pm$ SEM.

\section{Results}
\subsection{L7-Fmr1 KO mice are selectively impaired on LTD-dependent oculomotor learning}
VOR-increase learning: WT showed a significant increase across 30 min ($p = 7.37\times10^{-4}$); L7-Fmr1 KO did not ($p = 0.97$, 0 vs 30 min; L7-Fmr1 KO vs WT at 30 min $p = 3.47\times10^{-5}$). VOR-decrease learning was normal in L7-Fmr1 KO ($p = 1.10\times10^{-5}$, 0 vs 30 min; L7-Fmr1 KO vs WT at 30 min $p = 0.091$). The VOR-increase deficit was observed in both sexes (males, L7-Fmr1 KO vs WT at 30 min $p = 0.01$; females $p = 0.05$). Baseline VOR gain did not differ ($p = 0.95$, pre-training; $p = 0.50$ during paired visual+vestibular for VOR-increase; $p = 0.76$ for VOR-decrease).

\begin{table}[h]
\centering
\caption{Baseline VOR learning comparisons between L7-Fmr1 KO and WT mice.}
\begin{tabular}{lll}
\toprule
Comparison & Result \\
\midrule
VOR-increase, WT 0 vs 30 min & $p = 7.37\times10^{-4}$ (Tukey) \\
VOR-increase, L7-Fmr1 KO 0 vs 30 min & $p = 0.97$ (Tukey) \\
VOR-increase, L7-Fmr1 KO vs WT at 30 min & $p = 3.47\times10^{-5}$ (Tukey) \\
VOR-decrease, WT 0 vs 30 min & $p = 0.001$ (Tukey) \\
VOR-decrease, L7-Fmr1 KO 0 vs 30 min & $p = 1.10\times10^{-5}$ (Tukey) \\
VOR-decrease, L7-Fmr1 KO vs WT at 30 min & $p = 0.091$ (Tukey) \\
\bottomrule
\end{tabular}
\end{table}

\subsection{Behavioral pre-training rescues VOR-increase learning in L7-Fmr1 KO}
After 30 min of VOR-decrease pre-training, L7-Fmr1 KO mice showed VOR-increase learning indistinguishable from WT (VOR-decrease, L7-Fmr1 KO vs WT at 0 min: $p = 0.18$; subsequent VOR-increase, 0 vs 30 min: WT $p = 0.001$, L7-Fmr1 KO $p = 0.001$; L7-Fmr1 KO vs WT bar: $p = 0.17$; WT 30-min change with vs without pre-training: $p = 0.41$, paired sample t-test; L7-Fmr1 KO with vs without pre-training: $p = 0.01$, paired sample t-test). Vestibular-only pre-training produced similar effects (L7-Fmr1 KO vs WT VOR-increase 0 to 30 min: $p = 0.84$, paired sample t-test). In subsamples matched for mean VOR decrease during pre-training, VOR-increase was not significantly different between L7-Fmr1 KO and WT after VOR-decrease pre-training ($p = 0.74$) or Vestibular-only pre-training ($p = 0.40$).

\subsection{Diazepam pre-treatment (18--24 h) rescues VOR-increase and OKR deficits}
Diazepam did not alter baseline VOR gain (L7-Fmr1 KO: 2 h post $p = 0.72$, 18--24 h post $p = 0.77$; WT: 2 h post $p = 0.99$, 18--24 h post $p = 0.36$). The acute effect of diazepam (tested 2 h post injection, at both 0.4--0.5 mg/kg and 2.5 mg/kg) was to abolish VOR-increase learning in WT and L7-Fmr1 KO alike.

When VOR-increase training was delivered 18--24 h after diazepam (0.5 mg/kg), L7-Fmr1 KO mice showed robust learning comparable to WT ($p = 0.86$ L7-Fmr1 KO vs WT at 30 min; same individual L7-Fmr1 KO had no learning without pre-treatment, $p = 0.99$, 0 vs 30 min). The effect was transient: 1 week later, L7-Fmr1 KO again failed to learn ($p = 0.12$, 0 vs 30 min). Diazepam had no effect on WT VOR-increase learning ($p = 0.55$; same mice 1 week later $p = 0.99$), nor on VOR-decrease learning in either genotype (WT $p = 0.91$; L7-Fmr1 KO $p = 0.37$; L7-Fmr1 KO vs WT 1-day post-diazepam VOR-decrease at 30 min $p = 0.11$).

OKR adaptation: in WT, 60 min of training induced a learned increase ($p = 0.001$, 0 vs 60 min). L7-Fmr1 KO were significantly impaired (L7-Fmr1 KO vs WT at 60 min $p = 1.27\times10^{-4}$; L7-Fmr1 KO 0 vs 60 min $p = 0.87$). Diazepam pre-treatment 18--24 h before training fully rescued and even enhanced OKR learning in L7-Fmr1 KO (L7-Fmr1 KO saline vs diazepam at 60 min $p = 0.0001$; L7-Fmr1 KO vs WT post-diazepam $p = 0.02$ in favor of KO). Baseline OKR gain was unaffected (L7-Fmr1 KO vs WT: $p = 0.690$; L7-Fmr1 KO saline vs diazepam: $p = 0.690$; WT: $p = 0.55$).

\begin{table}[h]
\centering
\caption{Diazepam pre-treatment effects on LTD-dependent oculomotor learning (18--24 h delay).}
\begin{tabular}{lll}
\toprule
Task & Post-diazepam comparison & Result \\
\midrule
VOR-increase & L7-Fmr1 KO vs WT at 30 min & $p = 0.86$ (Tukey) \\
VOR-increase, 1 week & L7-Fmr1 KO 0 vs 30 min & $p = 0.12$ (Tukey) \\
VOR-decrease & L7-Fmr1 KO vs WT at 30 min & $p = 0.11$ (Tukey) \\
OKR adaptation & L7-Fmr1 KO vs WT at 60 min (diazepam) & $p = 0.02$ (KO $>$ WT, Tukey) \\
OKR adaptation & L7-Fmr1 KO saline vs diazepam at 60 min & $p = 0.0001$ (Tukey) \\
Males, VOR-increase & L7-Fmr1 KO vs WT at 30 min & $p = 0.290$ (Tukey) \\
Females, VOR-increase & L7-Fmr1 KO vs WT at 30 min & $p = 0.31$ (Tukey) \\
Males, OKR & L7-Fmr1 KO vs WT at 30 min & $p = 0.20$ (Tukey) \\
Females, OKR & L7-Fmr1 KO vs WT at 30 min & $p = 0.38$ (Tukey) \\
\bottomrule
\end{tabular}
\end{table}

\subsection{Low-frequency (0.5 Hz) VOR-increase impairment is not rescued}
Low-frequency VOR-increase learning showed a trend toward impairment in L7-Fmr1 KO ($p = 0.06$ L7-Fmr1 KO vs WT at 30 min). Neither VOR-decrease pre-training ($p = 0.47$, paired sample t-test), Vestibular-only pre-training ($p = 0.35$), nor 0.5 mg/kg diazepam ($p = 0.66$) rescued the low-frequency deficit, in contrast to their effectiveness at high frequency.

\begin{table}[h]
\centering
\caption{Task/frequency specificity of pre-training/pharmacological rescue.}
\begin{tabular}{lll}
\toprule
Task & L7-Fmr1 KO deficit & Rescued by pre-training / diazepam? \\
\midrule
High-freq (1 Hz) VOR-increase & yes ($p = 3.47\times10^{-5}$ vs WT) & yes ($p = 0.86$ vs WT post-diazepam) \\
OKR adaptation & yes ($p = 1.27\times10^{-4}$ vs WT) & yes ($p = 0.02$ vs WT, KO$>$WT) \\
VOR-decrease & no ($p = 0.091$) & not applicable \\
Low-freq (0.5 Hz) VOR-increase & trend ($p = 0.06$) & not rescued ($p = 0.47$, $p = 0.35$, $p = 0.66$) \\
\bottomrule
\end{tabular}
\end{table}

\section{Discussion}
In a Fragile X mouse model with enhanced PF-Purkinje cell LTD, learning deficits appeared selectively on LTD-dependent oculomotor tasks (high-frequency VOR-increase and OKR adaptation). These deficits were reversible by either behavioral pre-training designed to reverse LTD or by systemic diazepam pre-treatment 18--24 hours before training---a regimen that suppressed neural activity but allowed recovery of basal function. Together with prior results in MHCI K$^b$D$^b-/-$ mice \cite{nguyenvu2017saturation}, the findings support a sliding-threshold metaplasticity \cite{bienenstock1982theory,leet2022molecular} account: enhanced LTD is aberrantly recruited by ongoing activity, raising the threshold for further LTD and making the mechanism unavailable for learning.

The task specificity of both the deficit and its rescue is striking. High-frequency VOR-increase and OKR adaptation are the oculomotor tasks most strongly dependent on PF-Purkinje cell LTD \cite{boyden2006selective,hansel2006alphaca,guo2014biphasic,kimpo2014learning,kakegawa2018dynamics,jang2020limbic,shim2022ltp,zhang2023vestibular}. VOR-decrease learning was normal and unaffected, consistent with its independent mechanism (likely LTP of PF-Purkinje cell synapses \cite{shim2022ltp,levram2003synaptic}). Low-frequency VOR-increase learning was impaired but not rescuable, suggesting an additional, non-LTD cellular deficit from Fmr1 loss (e.g., mGluR-dependent dendritic translation \cite{huber2002altered}).

Diazepam's effectiveness suggests a clinical avenue: pharmacological suppression of neural activity during a pre-training period could reset elevated plasticity thresholds broadly, potentially addressing learning deficits in Fragile X syndrome and other conditions with enhanced associative plasticity. Pre-treatment did not impair baseline VOR, VOR-decrease learning, or baseline OKR, highlighting the specificity. Diazepam is FDA-approved and widely used, simplifying translation. This pharmacological approach contrasts with task-specific behavioral pre-training, which would require bespoke intervention per cerebellar function.

The threshold metaplasticity perspective extends the BCM framework \cite{bienenstock1982theory,abraham1996metaplasticity} to anti-Hebbian (LTD) synapses, where it cannot serve firing-rate homeostasis but may protect recently acquired memories and separate memories acquired in close succession onto different Purkinje cells \cite{fusi2005cascade,benna2016computational}. Related work has shown suppression of neural activity in the retina can reverse adult amblyopia \cite{fong2021reversal}, suggesting this may be a general strategy for resetting plasticity thresholds. Our findings argue that learning deficits from abnormal plasticity need not be permanent and can be remedied by appropriate circuit reset, opening therapeutic avenues.

\bibliographystyle{plain}
\bibliography{references}

\end{document}
